\chapter{Nonlinear Filtering: EKF, UKF and Particle Filters}
\label{ch:ch19}
% Function names for pseudocode are prefixed with "Nlf" (nonlinear filter)
% to stay unique across the book.
\SetKwFunction{NlfEkfStep}{EkfStep}
\SetKwFunction{NlfSigmaPoints}{SigmaPoints}
\SetKwFunction{NlfUnscentedTransform}{UnscentedTransform}
\SetKwFunction{NlfUkfStep}{UkfStep}
\SetKwFunction{NlfSystematicResample}{SystematicResample}
\SetKwFunction{NlfPfStep}{ParticleFilterStep}
\SetKwFunction{NlfWrap}{wrap}

\begin{objectives}
\begin{itemize}
  \item Explain why a turning intruder and a range--bearing sensor break the linear Kalman filter of \cref{ch:ch18}, and write down the coordinated-turn and range--bearing models with their Jacobians.
  \item Run one step of the extended Kalman filter (EKF) by hand: linearise, predict, wrap the bearing innovation, update; and say where the linearisation error comes from.
  \item Construct the $2n+1$ sigma points and weights of the unscented transform, propagate them through a nonlinear function, and recombine mean, covariance and cross-covariance into the unscented Kalman filter (UKF).
  \item Implement the bootstrap particle filter with systematic resampling, compute the effective sample size, and recognise degeneracy and sample impoverishment.
  \item Choose between EKF, UKF and particle filter for a given intruder model and sensor, and hand their output, including a multimodal particle cloud, to the prediction-aware planner of \cref{ch:ch24}.
\end{itemize}
\end{objectives}

% ---------------------------------------------------------------------
\section{Why this matters for a drone swarm}
\label{sec:ch19-motivation}

The Kalman filter of \cref{ch:ch18} gives the prediction layer of the hybrid
architecture (\cref{ch:ch24}) a stable estimate of an intruder's state and a
covariance that says how much to trust it. It rests on two assumptions: the
intruder moves according to a \emph{linear} model, and the sensor reports a
\emph{linear} function of the state. Both fail in the situations that matter
most. An unknown quadrotor that enters a survey site flies a straight leg,
banks into a turn around a crane and straightens out; a constant-velocity
model lags in every turn, and the natural model for a banked drone, the
\textbf{coordinated-turn model}\index{coordinated-turn model} with speed,
heading and turn rate in the state, contains the sine and cosine of the
heading. The sensor our drone carries, a radar, a lidar or a camera with a
rangefinder, reports how \emph{far} the intruder is and in which
\emph{direction}: a range and a bearing, the square root and the arctangent
of the position. The Kalman filter has no place for a sine or an arctangent.

This chapter presents the three standard answers in increasing order of
generality and cost. The \textbf{extended Kalman filter}\index{extended Kalman filter}
(EKF) linearises the models around the current estimate and runs the Kalman
equations on the linearisation. The \textbf{unscented Kalman filter}
\index{unscented Kalman filter} (UKF) avoids derivatives: it pushes a handful
of chosen points through the true functions and fits a Gaussian to what comes
out. The \textbf{particle filter}\index{particle filter} (PF) drops the
Gaussian assumption and represents the belief by thousands of weighted
samples, which lets it keep ``went left around the building'' and ``went
right'' alive at once. All three live in the prediction layer of
\cref{ch:ch24}: they consume raw measurements and deliver a state estimate
with its uncertainty to the local avoidance layer (\cref{ch:ch13,ch:ch14}),
to the trajectory predictor of \cref{ch:ch20} and to the prediction-aware
constraints of the planner; in the study plan (\cref{ch:appA}) they complete
Week~9. \Cref{sec:ch19-intuition,sec:ch19-problem} give the idea and the
models, \cref{sec:ch19-ekf,sec:ch19-ukf,sec:ch19-pf} develop the three
filters with pseudocode and worked examples and compare them on a turning
target, and
\cref{sec:ch19-properties,sec:ch19-variants,sec:ch19-implementation,sec:ch19-drone}
cover properties, variants, implementation and the place in the drone system.

% ---------------------------------------------------------------------
\section{The idea in plain words}
\label{sec:ch19-intuition}

Every filter step takes a Gaussian belief, pushes it through a function and
describes the result by a mean and a covariance. As long as the function is
affine, $\vect{y} = \mat{A}\state + \vect{b}$, the result is Gaussian with
mean $\mat{A}\vect{\mu} + \vect{b}$ and covariance $\mat{A}\mat{\Sigma}\mat{A}\T$
(\cref{ch:ch02}), and the Kalman filter is exact. As soon as the function
bends, the output is no longer Gaussian: push a Gaussian in range and bearing
through the polar-to-Cartesian map and you get a \emph{banana}, a curved
region that no ellipse describes well (\cref{fig:ch19-idea}). The three
filters are three approximations of that banana. The EKF replaces the bent
function by its tangent at the current estimate; it costs one Jacobian per
step, and the more the function curves within the width of the belief, the
larger the error. The UKF keeps the true function and changes the
representation of the input: $2n+1$ \emph{sigma points} whose sample mean and
covariance are exactly those of the belief are mapped through the true
function, and mean and covariance are read off the images; no derivative is
needed, and the mean is correct up to the second-order Taylor term, precisely
the term the EKF drops. The particle filter represents the belief by many
random samples with weights, moves every sample through the true function
with its own noise and re-weights the samples by how well they explain the
measurement; nothing forces the result to be an ellipse, so the belief can be
curved, skewed or split in two, and the price is the number of samples.
\index{linearisation error}\index{banana distribution}

\begin{figure}[tb]
  \centering
  \inputfigure{ch19/idea}
  \caption{Three ways of pushing a Gaussian belief (left) through a
  nonlinear function $f$ into the curved true output distribution (shaded
  banana). (a) The EKF moves the ellipse with the tangent at the mean; the
  ellipse sits on the image of the mean, not on the mean of the banana, and
  its ends miss the curved tails. (b) The UKF pushes $2n+1$ sigma points
  through the true $f$ and fits an ellipse to the images; the fitted mean
  lies inside the banana. (c) The particle filter pushes a cloud of random
  samples through $f$ and keeps the cloud as the belief, curvature included.}
  \label{fig:ch19-idea}
\end{figure}

\begin{keyidea}
A nonlinear filter has to push a belief through a bent function. The EKF
bends the belief with the tangent (one Jacobian), the UKF pushes $2n+1$
representative points through the real function (no Jacobian, second-order
accurate mean), and the particle filter pushes thousands of random samples
through it (no Gaussian at all, but $\bigO{N}$ work and the need to
resample). Choose the cheapest one whose assumptions your intruder and your
sensor satisfy.
\end{keyidea}

% ---------------------------------------------------------------------
\section{Problem statement and notation}
\label{sec:ch19-problem}

\begin{definition}[Nonlinear state-space model]
\label{def:ch19-model}
A \textbf{nonlinear state-space model}\index{state-space model!nonlinear} has
a state $\state_k \in \R^n$ at times $k = 0, 1, \dots$ spaced by $\dt$,
measurements $\meas_k \in \R^m$, a motion model $f$ and a measurement model
$h$ with
\begin{equation}
  \state_k = f(\state_{k-1}) + \vect{w}_k, \qquad
  \meas_k = h(\state_k) + \vect{v}_k,
  \label{eq:ch19-model}
\end{equation}
where $\vect{w}_k \sim \Normal(\vect{0}, \mat{Q}_k)$ and $\vect{v}_k \sim
\Normal(\vect{0}, \mat{R}_k)$ are independent of each other and over time, and
$\state_0 \sim \Normal(\hat{\state}_0, \mat{P}_0)$. The \textbf{filtering
problem}\index{filtering problem} is to compute, at every $k$, the posterior
$p(\state_k \mid \meas_{1:k})$ or at least its mean $\hat{\state}_k$ and
covariance $\mat{P}_k$.
\end{definition}

For $f(\state) = \mat{F}\state$ and $h(\state) = \mat{H}\state$ this is the
linear-Gaussian model of \cref{ch:ch18}, whose notation we keep:
$\hat{\state}_k^-$, $\mat{P}_k^-$ are the \emph{predicted} mean and covariance
before $\meas_k$ is used, $\hat{\state}_k$, $\mat{P}_k$ the \emph{updated}
ones, $\vect{y}_k$ the innovation, $\mat{S}_k$ its covariance and $\Kgain_k$
the gain. All three filters approximate the \textbf{Bayes filter}
\index{Bayes filter} \cite{thrun2005probabilistic},
\begin{equation}
  p(\state_k \mid \meas_{1:k-1}) = \!\int\! p(\state_k \mid \state_{k-1})\,
    p(\state_{k-1} \mid \meas_{1:k-1})\, \mathrm{d}\state_{k-1},
  \qquad
  p(\state_k \mid \meas_{1:k}) \propto p(\meas_k \mid \state_k)\,
    p(\state_k \mid \meas_{1:k-1}),
  \label{eq:ch19-bayes}
\end{equation}
which pushes yesterday's posterior through the motion model and then applies
Bayes' rule (\cref{ch:ch02}) with the likelihood $p(\meas_k \mid \state_k) =
\Normal(\meas_k; h(\state_k), \mat{R}_k)$. For linear-Gaussian models both
steps have closed forms; for the models below they do not, and each filter
replaces them by an approximation: a linearisation, a set of sigma points or
a set of samples. We write $\partial g/\partial\state$ for the
\textbf{Jacobian}\index{Jacobian} of $g$, the matrix of partial derivatives
$\partial g_i/\partial x_j$, which gives the best affine approximation
$g(\state + \vect{d}) \approx g(\state) + \frac{\partial g}{\partial
\state}(\state)\,\vect{d}$.

\subsection{The coordinated-turn model}
\label{sec:ch19-ct}

A drone flying at constant speed with a constant bank angle follows a
circular arc. The \textbf{coordinated-turn model}
\index{coordinated-turn model!definition} has the state
\begin{equation}
  \state = (x, y, v, \psi, \omega)\T,
  \label{eq:ch19-ct-state}
\end{equation}
with position $(x, y)$, speed $v \ge 0$, heading $\psi$ (counter-clockwise
from the $x$ axis) and turn rate $\omega = \dot{\psi}$; $\omega > 0$ turns
left, $\omega = 0$ is straight flight. Integrating $\dot{x} = v\cos\psi$,
$\dot{y} = v\sin\psi$, $\dot{\psi} = \omega$ over one step gives the exact
discrete model
\begin{equation}
  f(\state) =
  \begin{pmatrix}
    x + \dfrac{v}{\omega}\bigl(\sin(\psi + \omega\dt) - \sin\psi\bigr)\\[1.2ex]
    y - \dfrac{v}{\omega}\bigl(\cos(\psi + \omega\dt) - \cos\psi\bigr)\\[1.2ex]
    v\\
    \psi + \omega\dt\\
    \omega
  \end{pmatrix}
  \ (\omega \ne 0), \qquad
  f(\state) =
  \begin{pmatrix}
    x + v\dt\cos\psi\\
    y + v\dt\sin\psi\\
    v\\
    \psi\\
    0
  \end{pmatrix}
  \ (\omega = 0),
  \label{eq:ch19-ct}
\end{equation}
an arc of radius $v/\abs{\omega}$; the second form is the limit of the first,
used below a small threshold on $\abs{\omega}$. The model is nonlinear
through $v\sin\psi$, $v\cos\psi$ and the division by $\omega$. Bar-Shalom, Li
and Kirubarajan~\cite{barshalom2001estimation} write the same motion with
Cartesian velocities, $(x, \dot{x}, y, \dot{y}, \omega)$; the polar form has
the more compact Jacobian and delivers the speed and heading the avoidance
layer wants. Real intruders change speed and turn rate. We model a random
longitudinal acceleration $a_k \sim \Normal(0, \sigma_a^2)$ and a random turn
acceleration $\gamma_k \sim \Normal(0, \sigma_\gamma^2)$, constant during the
step, so that
\begin{equation}
  \vect{w}_k = \mat{G}(\state)\begin{pmatrix} a_k\\ \gamma_k\end{pmatrix},
  \quad
  \mat{G}(\state) =
  \begin{pmatrix}
    \tfrac12\dt^2\cos\psi & 0\\
    \tfrac12\dt^2\sin\psi & 0\\
    \dt & 0\\
    0 & \tfrac12\dt^2\\
    0 & \dt
  \end{pmatrix},
  \quad
  \mat{Q}(\state) = \mat{G}(\state)
  \begin{pmatrix}\sigma_a^2 & 0\\ 0 & \sigma_\gamma^2\end{pmatrix}
  \mat{G}(\state)\T,
  \label{eq:ch19-ct-noise}
\end{equation}
the white-noise-acceleration construction of \cref{ch:ch18} applied along
the heading and to the turn rate. $\mat{Q}$ depends on $\psi$, so the filters
evaluate it at the current estimate; $\sigma_a$ and $\sigma_\gamma$ say how
abruptly the intruder may manoeuvre.\index{process noise!coordinated turn}

\subsection{The range--bearing sensor and angles}
\label{sec:ch19-rb}

\begin{definition}[Range--bearing measurement]
\label{def:ch19-rb}
A sensor at the known position $\vect{s} = (s_x, s_y)$ measures the
\textbf{range}\index{range--bearing measurement} $r$ and the \textbf{bearing}
$\varphi \in (-\pi, \pi]$ of the intruder,
\begin{equation}
  h(\state) =
  \begin{pmatrix} r \\ \varphi \end{pmatrix} =
  \begin{pmatrix}
    \sqrt{(x - s_x)^2 + (y - s_y)^2}\\
    \operatorname{atan2}(y - s_y,\, x - s_x)
  \end{pmatrix},
  \qquad
  \vect{v}_k \sim \Normal\!\bigl(\vect{0}, \diag(\sigma_r^2, \sigma_\varphi^2)\bigr).
  \label{eq:ch19-rb}
\end{equation}
$h$ depends only on the position; the filter learns $v$, $\psi$ and $\omega$
from how the position changes over time.
\end{definition}

A bearing error of $\sigma_\varphi$ radians at range $r$ displaces the
position by about $r\sigma_\varphi$ across the line of sight, $5\,\mathrm{m}$
for $2^\circ$ at $150\,\mathrm{m}$; the error region is a thin arc, the
banana of \cref{fig:ch19-idea}. Headings and bearings live on a circle: the
difference between the bearings $3.10$ and $-3.10$ is not $6.20$ but
$-0.083$. We write
\begin{equation}
  \operatorname{wrap}(\theta) = \theta - 2\pi\Bigl\lceil \frac{\theta - \pi}{2\pi} \Bigr\rceil
  \in (-\pi, \pi]
  \label{eq:ch19-wrap}
\end{equation}
for the representative of $\theta$ on the circle; every difference of two
angles in this chapter, in an innovation, a residual or a likelihood, is
$\operatorname{wrap}(\theta_1 - \theta_2)$, and $\meas \ominus \hat{\meas}$
denotes a componentwise difference whose angular components are wrapped.
The heading of an estimate is wrapped after every update.\index{angle wrapping}

% ---------------------------------------------------------------------
\section{The extended Kalman filter}
\label{sec:ch19-ekf}

\subsection{Linearise, then run the Kalman filter}

The EKF replaces $f$ by its tangent at the current estimate and $h$ by its
tangent at the prediction,
\begin{equation}
  f(\state) \approx f(\hat{\state}_{k-1}) + \mat{F}_k(\state - \hat{\state}_{k-1}),
  \quad
  h(\state) \approx h(\hat{\state}_k^-) + \mat{H}_k(\state - \hat{\state}_k^-),
  \quad
  \mat{F}_k = \frac{\partial f}{\partial \state}(\hat{\state}_{k-1}),\
  \mat{H}_k = \frac{\partial h}{\partial \state}(\hat{\state}_k^-),
  \label{eq:ch19-ekf-linearise}
\end{equation}
and applies the Kalman equations of \cref{ch:ch18} to the tangents, with the
true $f$ and $h$ for the means:
\begin{align}
  \hat{\state}_k^- &= f(\hat{\state}_{k-1}),
  & \mat{P}_k^- &= \mat{F}_k \mat{P}_{k-1} \mat{F}_k\T + \mat{Q}_k,
  \label{eq:ch19-ekf-predict}\\
  \vect{y}_k &= \meas_k \ominus h(\hat{\state}_k^-),
  & \mat{S}_k &= \mat{H}_k \mat{P}_k^- \mat{H}_k\T + \mat{R}_k,
  \label{eq:ch19-ekf-innovation}\\
  \Kgain_k &= \mat{P}_k^- \mat{H}_k\T \mat{S}_k\inv,
  & \hat{\state}_k &= \hat{\state}_k^- + \Kgain_k \vect{y}_k,
  \qquad
  \mat{P}_k = (\mat{I} - \Kgain_k \mat{H}_k)\,\mat{P}_k^-.
  \label{eq:ch19-ekf-update}
\end{align}
Only the covariances and the gain use the linearisation. The Jacobians are
evaluated at the current estimate, so they change from step to step and
depend on the data: the EKF covariance cannot be computed in advance, and a
wrong estimate produces a wrong Jacobian and a wrong gain
(\cref{sec:ch19-ekf-failure}).\index{extended Kalman filter!equations}

\begin{algorithm}[htb]
\caption{One step of the extended Kalman filter}
\label{alg:ch19-ekf}
\KwIn{previous estimate $\hat{\state}_{k-1}$, $\mat{P}_{k-1}$; measurement $\meas_k$ or \KwNil; models $f$, $h$ with Jacobians; $\mat{Q}_k$, $\mat{R}_k$}
\KwOut{updated estimate $\hat{\state}_k$, $\mat{P}_k$}
\Fn{\NlfEkfStep{$\hat{\state}_{k-1}$, $\mat{P}_{k-1}$, $\meas_k$}}{
  $\mat{F}_k \gets \partial f / \partial\state$ at $\hat{\state}_{k-1}$; $\hat{\state}_k^- \gets f(\hat{\state}_{k-1})$, heading wrapped; $\mat{P}_k^- \gets \mat{F}_k \mat{P}_{k-1} \mat{F}_k\T + \mat{Q}_k$\label{alg:ch19-ekf:predict}\;
  \If{$\meas_k = $ \KwNil}{\KwRet $(\hat{\state}_k^-, \mat{P}_k^-)$\tcp*{missing measurement: predict only}\label{alg:ch19-ekf:missing}}
  $\mat{H}_k \gets \partial h / \partial\state$ at $\hat{\state}_k^-$\label{alg:ch19-ekf:H}\;
  $\vect{y}_k \gets \meas_k - h(\hat{\state}_k^-)$; $\vect{y}_k[j] \gets$ \NlfWrap{$\vect{y}_k[j]$} for every angular component $j$\label{alg:ch19-ekf:innovation}\;
  $\mat{S}_k \gets \mat{H}_k \mat{P}_k^- \mat{H}_k\T + \mat{R}_k$; $\Kgain_k \gets \mat{P}_k^- \mat{H}_k\T \mat{S}_k\inv$\label{alg:ch19-ekf:gain}\;
  $\hat{\state}_k \gets \hat{\state}_k^- + \Kgain_k \vect{y}_k$, heading wrapped\label{alg:ch19-ekf:update}\;
  $\mat{P}_k \gets (\mat{I} - \Kgain_k \mat{H}_k)\,\mat{P}_k^- (\mat{I} - \Kgain_k \mat{H}_k)\T + \Kgain_k \mat{R}_k \Kgain_k\T$\tcp*{Joseph form}\label{alg:ch19-ekf:joseph}
  \KwRet $(\hat{\state}_k, \mat{P}_k)$\;
}
\end{algorithm}

Line~\ref{alg:ch19-ekf:predict} of \cref{alg:ch19-ekf} is the predict step,
with $\mat{F}_k$ at the \emph{previous} estimate; line~\ref{alg:ch19-ekf:missing}
handles a step without measurement, as when an obstacle hides the intruder,
by taking the prediction as the estimate. The update linearises $h$ at the
\emph{predicted} estimate (line~\ref{alg:ch19-ekf:H}); the innovation on
line~\ref{alg:ch19-ekf:innovation} has its bearing wrapped, without which an
intruder near the $\pm\pi$ seam produces an innovation of almost $2\pi$
(\cref{sec:ch19-ekf-example}). Line~\ref{alg:ch19-ekf:joseph} uses the
Joseph form of the covariance update, algebraically equal to $(\mat{I} -
\Kgain_k\mat{H}_k)\mat{P}_k^-$ but symmetric and positive definite under
rounding.

\subsection{The Jacobians of the two models}
\label{sec:ch19-ekf-jacobians}

\begin{proposition}[Jacobian of the coordinated-turn model]
\label{thm:ch19-ct-jacobian}
Let $\psi' = \psi + \omega\dt$. For $\omega \ne 0$ the Jacobian of
\cref{eq:ch19-ct} is
\begin{equation}
  \mat{F} =
  \begin{pmatrix}
    1 & 0 & \dfrac{\sin\psi' - \sin\psi}{\omega}
      & \dfrac{v}{\omega}\bigl(\cos\psi' - \cos\psi\bigr)
      & \dfrac{v\dt\cos\psi'}{\omega} - \dfrac{v\,(\sin\psi' - \sin\psi)}{\omega^2}\\[1.6ex]
    0 & 1 & -\dfrac{\cos\psi' - \cos\psi}{\omega}
      & \dfrac{v}{\omega}\bigl(\sin\psi' - \sin\psi\bigr)
      & \dfrac{v\dt\sin\psi'}{\omega} + \dfrac{v\,(\cos\psi' - \cos\psi)}{\omega^2}\\[1.6ex]
    0 & 0 & 1 & 0 & 0\\
    0 & 0 & 0 & 1 & \dt\\
    0 & 0 & 0 & 0 & 1
  \end{pmatrix},
  \label{eq:ch19-ct-jacobian}
\end{equation}
and for $\omega = 0$ the first two rows become
$F_{13} = \dt\cos\psi$, $F_{14} = -v\dt\sin\psi$,
$F_{15} = -\tfrac12 v\dt^2\sin\psi$, $F_{23} = \dt\sin\psi$,
$F_{24} = v\dt\cos\psi$, $F_{25} = \tfrac12 v\dt^2\cos\psi$.
\index{Jacobian!coordinated turn}
\end{proposition}
\begin{proof}[Proof sketch]
Rows three to five are the derivatives of $v$, $\psi + \omega\dt$ and
$\omega$. In the first row, $v$ enters linearly with the factor
$(\sin\psi' - \sin\psi)/\omega$; differentiating with respect to $\psi$ uses
$\partial\psi'/\partial\psi = 1$, and with respect to $\omega$ the quotient
rule on $1/\omega$ together with $\partial\psi'/\partial\omega = \dt$; the
second row is the same computation for $-\frac{v}{\omega}(\cos\psi' -
\cos\psi)$. The $\omega = 0$ entries are limits: $\sin\psi' = \sin\psi +
\omega\dt\cos\psi - \tfrac12\omega^2\dt^2\sin\psi + \bigO{\omega^3}$ gives
$x' = x + v\dt\cos\psi - \tfrac12 v\omega\dt^2\sin\psi + \bigO{\omega^2}$,
whose $\omega$-derivative at $0$ is $-\tfrac12 v\dt^2\sin\psi$.
\Cref{exr:ch19-ct-jacobian} completes the limits; the self-test of the
chapter code checks every entry against finite differences.
\end{proof}

\begin{proposition}[Jacobian of the range--bearing model]
\label{thm:ch19-rb-jacobian}
With $d_x = x - s_x$, $d_y = y - s_y$ and $r = \sqrt{d_x^2 + d_y^2} > 0$,
\begin{equation}
  \mat{H} = \frac{\partial h}{\partial \state} =
  \begin{pmatrix}
    d_x/r & d_y/r & 0 & 0 & 0\\[0.6ex]
    -d_y/r^2 & d_x/r^2 & 0 & 0 & 0
  \end{pmatrix}.
  \label{eq:ch19-rb-jacobian}
\end{equation}
\index{Jacobian!range--bearing}
\end{proposition}
\begin{proof}
The chain rule on the square root gives the first row; for the bearing,
$\varphi = \operatorname{atan2}(d_y, d_x)$ has $\partial\varphi/\partial x =
-d_y/(d_x^2 + d_y^2)$ and $\partial\varphi/\partial y = d_x/(d_x^2 + d_y^2)$
wherever $r > 0$, on every branch of the arctangent. The last three columns
vanish because $h$ ignores $v$, $\psi$, $\omega$.
\end{proof}

The first row of $\mat{H}$ is the unit vector from the sensor to the
intruder: a range constrains the position along the line of sight. The
second row is the perpendicular unit vector divided by $r$: a bearing
constrains the position across the line of sight, and the further the
intruder, the less a radian of bearing is worth.

\subsection{A worked example: one EKF step by hand}
\label{sec:ch19-ekf-example}

\begin{example}[One EKF step with a range--bearing measurement]
\label{ex:ch19-ekf}
Our drone hovers at $\vect{s} = (0, 0)$ and observes an intruder once per
second ($\dt = 1\,\mathrm{s}$). After step $k-1$,
\[
  \hat{\state}_{k-1} = (-100,\ -4,\ 10,\ \pi,\ 0.3)\T, \qquad
  \mat{P}_{k-1} = \diag(4,\ 4,\ 1,\ 0.04,\ 0.01),
\]
in metres, metres per second, radians and radians per second: the intruder
is $100\,\mathrm{m}$ to the west, flies west at $10\,\mathrm{m/s}$ and turns
left at $0.3\,\mathrm{rad/s}$, so its path curves towards the south. The
process noise has $\sigma_a = 1\,\mathrm{m/s^2}$ and $\sigma_\gamma =
0.1\,\mathrm{rad/s^2}$, the sensor $\sigma_r = 2\,\mathrm{m}$ and
$\sigma_\varphi = 0.03\,\mathrm{rad}$, and the new measurement is $\meas_k =
(111.5\,\mathrm{m},\ 3.12\,\mathrm{rad})$. The instance is
\code{worked\_example()} in \code{ch19\_nonlinear\_filters.py}; every number
below is printed by its self-test.
\end{example}

\paragraph{Predict.}
With $v/\omega = 33.33$ and $\psi' = \pi + 0.3$, \cref{eq:ch19-ct} gives
$x^- = -100 + 33.33\sin\psi' = -109.85$ and $y^- = -4 - 33.33\,(\cos\psi' + 1)
= -5.49$: $9.85\,\mathrm{m}$ west and $1.49\,\mathrm{m}$ south. The heading
$\pi + 0.3$ wraps to $-2.8416$; speed and turn rate are unchanged. The
Jacobian \cref{eq:ch19-ct-jacobian} at $\hat{\state}_{k-1}$ has $F_{13} =
-0.985$, $F_{14} = 1.489$, $F_{15} = 0.991$, $F_{23} = -0.149$, $F_{24} =
-9.851$ and $F_{25} = -4.888$: a heading error of $0.1\,\mathrm{rad}$ moves
the predicted position by $1\,\mathrm{m}$ across the direction of flight
($F_{24}$). The predicted covariance $\mat{F}_k\mat{P}_{k-1}\mat{F}_k\T +
\mat{Q}_k$ has the position block $\begin{psmallmatrix} 5.319 & -0.488\\
-0.488 & 8.143\end{psmallmatrix}$ (standard deviations $2.31$ and
$2.85\,\mathrm{m}$, larger across the westward flight) and now carries
covariances of $-0.443$ between $y$ and $\psi$ and $-1.485$ between $x$ and
$v$, which the update will use to correct heading and speed from a position
measurement.

\paragraph{Update.}
The predicted measurement is $h(\hat{\state}_k^-) = (109.99, -3.0917)$: the
intruder is expected just south of the negative $x$ axis, so its bearing is
just above $-\pi$, while the measured bearing $3.12$ is just below $+\pi$.
Subtracting naively gives $6.2117$, an innovation of $356^\circ$ that would
send the estimate to the other side of the sensor; wrapping gives
$-0.0715\,\mathrm{rad}$, so $\vect{y}_k = (1.51,\ -0.0715)$: $1.5\,\mathrm{m}$
further away and $4.1^\circ$ clockwise of the prediction. With
$\mat{H}_k = \begin{psmallmatrix} -0.9988 & -0.0499 & 0 & 0 & 0\\
0.00045 & -0.00908 & 0 & 0 & 0\end{psmallmatrix}$ the innovation covariance
is $\mat{S}_k \approx \diag(9.28, 0.0016)$, and the gain maps the range
innovation mostly into $x$ and the bearing innovation mostly into $y$:
$K_{11} = -0.569$ and $K_{22} = -47.1$. The second value is no typo: a radian
of bearing at $110\,\mathrm{m}$ is $110\,\mathrm{m}$ across the line of
sight, so the bearing gain is of the order of $-r$. The rows for the
unmeasured components come from the correlations in $\mat{P}_k^-$: $K_{42} =
2.56$ turns a bearing innovation into a heading correction, $K_{31} = 0.16$ a
range innovation into a speed correction. \Cref{tab:ch19-ekf-trace} lists the
result: the position moved $1.09\,\mathrm{m}$ west and $3.35\,\mathrm{m}$
north, almost entirely from the bearing; the heading turned by
$-0.19\,\mathrm{rad}$ towards due west (an intruder found north of the
predicted arc is best explained by a less southerly heading); the speed rose
by $0.19\,\mathrm{m/s}$. The updated position standard deviations are
$1.51\,\mathrm{m}$ along the line of sight, set by the range noise, and
$2.16\,\mathrm{m}$ across it, set by the bearing noise at this range.
\Cref{exr:ch19-wrapping} asks what the unwrapped innovation would have done.

\begin{table}[tb]
  \centering\small
  \caption{One step of \cref{alg:ch19-ekf} on \cref{ex:ch19-ekf} and the
  same step of the UKF (\cref{sec:ch19-ukf}). The two filters agree to about
  a decimetre because the heading uncertainty is small compared with the
  curvature of the arc; \cref{sec:ch19-ukf-compare} shows cases where they
  do not.}
  \label{tab:ch19-ekf-trace}
  \begin{tabular}{@{}lrrrrr@{}}
  \toprule
  Quantity & $x$ [m] & $y$ [m] & $v$ [m/s] & $\psi$ [rad] & $\omega$ [rad/s]\\
  \midrule
  $\hat{\state}_{k-1}$ & $-100.000$ & $-4.000$ & $10.000$ & $3.1416$ & $0.3000$\\
  EKF predicted $\hat{\state}_k^-$ & $-109.851$ & $-5.489$ & $10.000$ & $-2.8416$ & $0.3000$\\
  EKF updated $\hat{\state}_k$ & $-110.941$ & $-2.134$ & $10.190$ & $-3.0313$ & $0.2787$\\
  UKF predicted $\hat{\state}_k^-$ & $-109.641$ & $-5.456$ & $10.000$ & $-2.8416$ & $0.3000$\\
  UKF updated $\hat{\state}_k$ & $-110.834$ & $-2.176$ & $10.214$ & $-3.0284$ & $0.2784$\\
  \midrule
  std.\ dev.\ before & $2.000$ & $2.000$ & $1.000$ & $0.2000$ & $0.1000$\\
  std.\ dev.\ predicted (EKF) & $2.306$ & $2.854$ & $1.414$ & $0.2291$ & $0.1414$\\
  std.\ dev.\ updated (EKF) & $1.513$ & $2.157$ & $1.326$ & $0.2047$ & $0.1410$\\
  std.\ dev.\ predicted (UKF) & $2.355$ & $2.810$ & $1.414$ & $0.2291$ & $0.1414$\\
  std.\ dev.\ updated (UKF) & $1.527$ & $2.137$ & $1.329$ & $0.2059$ & $0.1410$\\
  \bottomrule
  \end{tabular}
\end{table}

\subsection{Where the EKF fails}
\label{sec:ch19-ekf-failure}

\paragraph{Strong nonlinearity.}
\Cref{fig:ch19-linearisation} makes the linearisation error visible. A
Gaussian belief in range and bearing, $r \sim \Normal(100, 5^2)$ and
$\varphi \sim \Normal(\pi/2, 0.3^2)$, is pushed through $(r, \varphi) \mapsto
(r\cos\varphi, r\sin\varphi)$, which is what a filter does when it
initialises a track from one measurement. The true output, estimated from
$2\,000$ samples, has its mean at $(-0.2, 95.6)$; the exact mean is
$100\,e^{-0.3^2/2} = 95.6$, because the samples lie on an arc of radius $100$
and the average of points on an arc is \emph{inside} the arc. The linearised
map puts the mean at the image of the input mean, $(0, 100)$,
$4.4\,\mathrm{m}$ too far out, and its ellipse is a straight bar that cannot
follow the curved tails: its standard deviation along $y$ is $5.0\,\mathrm{m}$
where the samples spread by $7.7\,\mathrm{m}$, and its $2\sigma$ ellipse
contains $74.6\,\%$ of the samples instead of the nominal $86.5\,\%$, missing
exactly the tails a collision check cares about. Nothing in the EKF signals
the error: it reports a confident covariance around a biased mean. The
unscented transform of \cref{sec:ch19-ukf} puts the mean at $(0.0, 95.6)$
with five function evaluations and no derivative, and its ellipse contains
$91.6\,\%$ of the samples.

\begin{figure}[tb]
  \centering
  \inputfigure{ch19/linearisation}
  \caption{The linearisation error of the EKF on the polar-to-Cartesian map.
  Grey dots: $2\,000$ samples of a Gaussian in $(r, \varphi)$ mapped to the
  plane; they form a banana. Left: the EKF's $2\sigma$ ellipse (blue) is
  centred on the image of the mean, $4.4\,\mathrm{m}$ beyond the true mean
  (cross), and cannot bend. Right: the five sigma points of the unscented
  transform mapped through the same function (orange squares) and the
  ellipse fitted to them, centred inside the banana. The dashed ellipse is
  the $2\sigma$ ellipse of the samples themselves.}
  \label{fig:ch19-linearisation}
\end{figure}

\paragraph{A poor initial estimate.}
If the estimate is far from the truth, $\mat{F}_k$ and $\mat{H}_k$ describe
the wrong tangent, the gain is computed for a fictitious geometry, and an
update can push the estimate further away while the covariance, which knows
nothing about the error, keeps shrinking: the filter has
\textbf{diverged}\index{divergence (EKF)}. Whether this happens depends on
how strongly the measurements constrain the state. A range--bearing sensor
measures the position directly, and in the experiment of
\cref{sec:ch19-ukf-compare} the EKF recovers from a $90^\circ$ heading error
in every run; divergence is the real danger with weak or bearing-only
measurements, long gaps, and an initial covariance small enough to make the
filter ignore its innovations. The cures are an honest $\mat{P}_0$, a filter
whose approximation does not depend on a tangent, and, when the measurements
cannot resolve the ambiguity, a particle filter that keeps several
hypotheses alive.

